\chapter{Giới thiệu}
\label{chap:introduction}

\section{Bối cảnh và động lực}

Mô hình ngôn ngữ định hướng suy luận thường cải thiện kết quả khi được cấp thêm
tài nguyên ở thời điểm suy luận. Hai cơ chế phổ biến là tăng giới hạn token cho
một quỹ đạo sinh và lấy nhiều mẫu độc lập rồi tổng hợp đáp án. Chain-of-thought
giúp mô hình biểu diễn các bước trung gian~\cite{wei2022cot}; self-consistency
tăng độ tin cậy bằng cách tổng hợp nhiều quỹ đạo~\cite{wang2023selfconsistency};
và các nghiên cứu về repeated sampling cho thấy xác suất tìm được lời giải đúng
có thể tăng khi số mẫu tăng~\cite{brown2024monkeys}. Những kết quả này tạo cơ
sở cho test-time scaling nhưng chưa trả lời trực tiếp cách phân bổ compute theo
từng bài trong điều kiện phần cứng giới hạn.

Một cấu hình duy nhất cho mọi đầu vào có hai nhược điểm. Thứ nhất, bài dễ có thể
được giải đúng từ mẫu đầu tiên nhưng vẫn phải trả chi phí cho các mẫu sau. Thứ
hai, nhiều quỹ đạo ngắn không thể bù cho một giới hạn token quá thấp nếu mọi
quỹ đạo đều bị cắt trước khi mô hình xuất đáp án. Vấn đề vì vậy không chỉ là
``dùng nhiều compute hơn'', mà là chọn hình thức compute phù hợp với trạng thái
quan sát được của từng bài.

Bài toán đặc biệt quan trọng với mô hình cỡ nhỏ chạy cục bộ. Mô hình nhỏ phù hợp
với GPU phổ thông nhưng có biên độ sai số lớn hơn, còn việc tăng đồng thời số
mẫu và giới hạn token nhanh chóng làm tăng thời gian chạy. Khóa luận tập trung
vào Qwen3-0.6B trên một GPU RTX 4060, qua đó đặt bài toán trong điều kiện có thể
tái lập ở quy mô khóa luận đại học thay vì giả định hạ tầng suy luận lớn.

\section{Phát biểu bài toán}

Với bài toán $x_i$ có đáp án chuẩn $y_i$, một cấu hình suy luận được ký hiệu
\begin{equation}
  p=(k,L),
\end{equation}
trong đó $k$ là số quỹ đạo sinh độc lập và $L$ là giới hạn
\code{max\_new\_tokens} của mỗi quỹ đạo. Khóa luận gọi $k$ là \emph{breadth}.
Từ \emph{depth} chỉ được dùng như cách viết tắt cho giới hạn token $L$, không
được diễn giải là độ sâu ngữ nghĩa của quá trình suy luận.

Ngân sách danh nghĩa của cấu hình là $kL$, nhưng chi phí thực tế được đo bằng
\begin{equation}
  C_i(p)=\sum_{j=1}^{k} t_{ij}(L), \qquad t_{ij}(L)\leq L,
\end{equation}
với $t_{ij}$ là số token thật sự đã sinh. Phân biệt này cần thiết vì mô hình có
thể kết thúc sớm hoặc chạm giới hạn ở từng quỹ đạo. Mọi mẫu mà chính sách đã yêu
cầu đều được tính chi phí, kể cả khi mẫu đó không quyết định đáp án cuối.

Khóa luận xem xét liệu một chính sách nhẹ, không dùng verifier ngoài và không
truy cập hidden state, có thể quyết định tăng $L$, lấy thêm mẫu, hoặc dừng dựa
trên thông tin đã thanh toán hay không. Chính sách không được dùng đáp án chuẩn,
nhãn oracle, mẫu tương lai hoặc kết quả phản thực tế tại thời điểm ra quyết
định.

\section{Câu hỏi nghiên cứu}

Câu hỏi nghiên cứu đã chốt của đề tài là:

\begin{quote}
Đối với mô hình suy luận mã nguồn mở cỡ nhỏ, liệu một chính sách tuần tự nhẹ,
không sử dụng verifier, có khả năng phân bổ giới hạn token sinh và chọn lọc bổ
sung các mẫu độc lập, có thể cải thiện biên Pareto giữa độ chính xác trên tập
giữ lại và số token sinh thực tế so với các cấu hình cố định cùng baseline
thích ứng chỉ theo depth hoặc breadth hay không?
\end{quote}

Câu hỏi này tương ứng với phát biểu tiếng Anh được dùng trong paper:

\begin{otherlanguage}{english}
\begin{quote}
For small open-weight reasoning models, can a verifier-free lightweight
sequential policy that allocates generation-token caps and selectively adds
independent samples improve the held-out accuracy--actual-generated-token
Pareto frontier over fixed profiles and depth-only or breadth-only adaptive
baselines?
\end{quote}
\end{otherlanguage}

Hai câu hỏi hỗ trợ được dùng để diễn giải kết quả: (1) khi nào selective
breadth mang lại giá trị ngoài adaptive token cap; và (2) quyết định phân bổ có
ổn định giữa các tập bài suy luận hay không. Đánh giá trên một họ mô hình thứ
hai vẫn là hướng mở và không được xem là phần bằng chứng đã hoàn tất.

\section{Mục tiêu nghiên cứu}

\subsection{Mục tiêu tổng quát}

Mục tiêu tổng quát là xây dựng và đánh giá một quy trình phân bổ test-time
compute có thể tái lập cho mô hình suy luận cỡ nhỏ, trong đó độ chính xác được
đánh đổi với chi phí token thực tế và quyết định thích ứng không sử dụng tín
hiệu bị rò rỉ từ đáp án chuẩn hoặc dữ liệu Test.

\subsection{Mục tiêu cụ thể}

\begin{enumerate}
  \item Xây dựng pipeline sinh, lưu cache, tiếp tục sau gián đoạn, parse đáp án,
  tổng hợp và phát lại chính sách trên cùng kho mẫu.
  \item Định lượng ảnh hưởng riêng và kết hợp của $k$ và $L$ trên GSM8K và
  MATH-500 bằng cấu hình, split và policy được đóng băng.
  \item So sánh chính sách cố định, depth-only, breadth-only và heuristic kết
  hợp theo độ chính xác và số token thực tế.
  \item Báo cáo độ bất định và so sánh ghép cặp, đồng thời tách kết quả Test mô
  tả khỏi quá trình chọn chính sách trên Validation.
  \item Xác định giới hạn của kết luận trong điều kiện một seed, một họ mô hình
  chính và phần cứng cục bộ.
\end{enumerate}

\section{Phạm vi và các quyết định phương pháp}

Mô hình chính được chốt là Qwen3-0.6B. Hai tập dữ liệu chính là GSM8K và
MATH-500, mỗi tập có 100 bài Test giữ lại, tách rời Train và Validation. GSM8K
dùng $L\in\{512,1024,2048\}$; MATH-500 dùng
$L\in\{1024,2048,4096\}$; cả hai dùng $k\in\{1,2,4\}$. Sinh mẫu sử dụng
thinking mode, temperature 0,7, top-$p$ 0,95 và seed 42. Phân tích lịch sử
Qwen3-1.7B/GSM8K được giữ thành một tầng bằng chứng riêng, không gộp với kết quả
0.6B.

Khóa luận không đặt mục tiêu huấn luyện lại mô hình nền, xây dựng verifier, hoặc
tuyên bố một controller học máy hoàn chỉnh. Dữ liệu Train không có đủ sự kiện
escalation dương để huấn luyện đáng tin cậy một joint controller. Vì vậy chính
sách kết hợp được đánh giá là heuristic tất định đã đóng băng. Qwen3-1.7B trên
MATH-500 và lưới $L=8192$ không thuộc điều kiện hoàn thành vì không phù hợp với
ngân sách thời gian và bộ nhớ GPU hiện có.

\section{Đóng góp đã hoàn thành}

Các đóng góp được giới hạn ở những thành phần đã có artifact kiểm chứng:

\begin{enumerate}
  \item Pipeline đánh giá breadth--token-cap có manifest, seed tất định, cache
  append-only, tiếp tục sau gián đoạn, parser nghiêm ngặt, cost accounting theo
  token thực tế và phát lại policy.
  \item Đánh giá giữ lại của 14 chính sách đã đăng ký trên Test-100 của GSM8K
  và MATH-500, không chọn lại chính sách dựa trên Test.
  \item Bằng chứng rằng dừng theo mức đồng thuận có thể giữ nguyên độ chính xác
  so với fixed high-breadth trong khi giảm 21,3\% token trên GSM8K và xấp xỉ
  20\% token trên MATH-500.
  \item Kết quả âm có giá trị: heuristic kết hợp không tạo cải thiện Pareto ổn
  định giữa hai tập dữ liệu và thất bại rõ rệt trên MATH-500.
\end{enumerate}

Non-thinking mode, native thinking budget, budget forcing, seed bổ sung và
controller học máy là phương pháp hoặc đánh giá tương lai. Chúng không được
liệt kê như đóng góp đã hoàn thành nếu chưa có kết quả được kiểm tra.

\section{Tổng quan phương pháp nghiên cứu}

Quy trình nghiên cứu gồm bốn giai đoạn. Giai đoạn đầu khảo sát lịch sử với
Qwen3-1.7B trên GSM8K để hình thành giả thuyết ``đảm bảo token cap trước, bổ sung
breadth có chọn lọc''. Giai đoạn hai chuyển sang Qwen3-0.6B, chuẩn hóa parser,
manifest và action space phù hợp với RTX 4060. Giai đoạn ba dùng Train và
Validation để xây dựng dữ liệu tuần tự, xác định policy, kiểm tra leakage và
đóng băng điểm vận hành. Giai đoạn bốn sinh đúng một lần trên hai tập Test-100,
phát lại toàn bộ policy và thực hiện phân tích ghép cặp.

Cách tổ chức này phân biệt ba loại phát biểu. \emph{Bằng chứng hoàn thành} phải
có raw artifact, manifest và phép phân tích tái lập. \emph{Giả thuyết làm việc}
dùng để giải thích hoặc đề xuất thí nghiệm nhưng không được viết như kết luận.
\emph{Đánh giá tương lai} là phần chưa chạy và chỉ xuất hiện trong giới hạn hoặc
hướng phát triển.

\section{Cấu trúc khóa luận}

Sau chương giới thiệu, Chương 2 trình bày nền tảng và công trình liên quan.
Chương 3 định nghĩa bài toán, chính sách và phương pháp thống kê. Chương 4 mô tả
kiến trúc phần mềm cùng cơ chế tái lập. Chương 5 trình bày thiết kế thực nghiệm
đã đóng băng. Chương 6 báo cáo kết quả GSM8K và MATH-500. Chương 7 thảo luận ý
nghĩa, đe dọa độ giá trị và giới hạn. Chương 8 trả lời câu hỏi nghiên cứu và nêu
hướng phát triển. Phụ lục lưu cấu hình, lệnh chạy, bảng đầy đủ và danh mục
artifact hỗ trợ kiểm chứng.
